
\documentclass{article}
\usepackage{colortbl}
\usepackage{makecell}
\usepackage{multirow}
\usepackage{supertabular}

\begin{document}

\newcounter{utterance}

\centering \large Interaction Transcript for game `cladder', experiment `full\_v1.5\_default', episode 4422 with refineonly{-}Qwen3.5{-}27B.
\vspace{24pt}

{ \footnotesize  \setcounter{utterance}{1}
\setlength{\tabcolsep}{0pt}
\begin{supertabular}{c@{$\;$}|p{.15\linewidth}@{}p{.15\linewidth}p{.15\linewidth}p{.15\linewidth}p{.15\linewidth}p{.15\linewidth}}
    \# & \multicolumn{2}{c}{Player} && \multicolumn{2}{c}{Game Master} \\
    \hline

    \theutterance \stepcounter{utterance}  
    & & & \multicolumn{4}{p{0.6\linewidth}}{
        \cellcolor[rgb]{0.9,0.9,0.9}{
            \makecell[{{p{\linewidth}}}]{
                \texttt{\tiny{[P1$\langle$GM]}}
                \texttt{You are an expert in causal inference. The following question is not a typical commonsense query, but rather a meticulously designed question created by a professor specializing in causal inference, intended to assess the students' mastery of the course content.} \\
\texttt{Imagine a self{-}contained, hypothetical world with only the following conditions, and without any unmentioned factors or causal relationships: Gender has a direct effect on smoking and peanut allergy. Smoking has a direct effect on tar deposit. Tar deposit has a direct effect on peanut allergy. Gender is unobserved. For nonsmokers, the probability of high tar deposit is 18\%. For smokers, the probability of high tar deposit is 64\%. For nonsmokers and with no tar deposit, the probability of being allergic to peanuts is 46\%. For nonsmokers and with high tar deposit, the probability of being allergic to peanuts is 77\%. For smokers and with no tar deposit, the probability of being allergic to peanuts is 17\%. For smokers and with high tar deposit, the probability of being allergic to peanuts is 43\%. The overall probability of smoking is 49\%. Will smoking decrease the chance of being allergic to peanuts?} \\
\texttt{Start your answer with 'Yes' or 'No', followed by additional reasoning or evidence to support your explanation.} \\
            }
        }
    }
     \\ \\

    \theutterance \stepcounter{utterance}  
    & \multicolumn{4}{p{0.6\linewidth}}{
        \cellcolor[rgb]{0.561,0.737,0.561}{
            \makecell[{{p{\linewidth}}}]{
                \texttt{\tiny{[P1$\rangle$GM]}}
                \texttt{No} \\
\\ 
\texttt{To determine whether smoking decreases the chance of being allergic to peanuts, we must calculate the **Average Treatment Effect (ATE)** of smoking on peanut allergy. This requires comparing the expected probability of allergy if everyone smoked versus if no one smoked.} \\
\\ 
\texttt{However, a critical complication exists: **Gender** is an unobserved confounder. It affects both Smoking (\$S\$) and Peanut Allergy (\$A\$). Because Gender is unobserved, we cannot simply stratify by it to remove confounding bias. Furthermore, the problem states that the probabilities of allergy are given conditional on *both* Smoking status and Tar Deposit status (\$P(A|S, T)\$), but the distribution of Tar Deposit (\$T\$) depends on Smoking (\$S\$).} \\
\\ 
\texttt{Let's break down the causal structure and the available data:} \\
\texttt{1.  **Causal Graph**: \$Gender \textbackslash{}to Smoking\$, \$Gender \textbackslash{}to Allergy\$, \$Smoking \textbackslash{}to Tar\$, \$Tar \textbackslash{}to Allergy\$.} \\
\texttt{2.  **Goal**: Calculate \$E{[}A|do(S=1){]} {-} E{[}A|do(S=0){]}\$.} \\
\texttt{3.  **Data Provided**:} \\
\texttt{    *   \$P(T=1|S=0) = 0.18\$} \\
\texttt{    *   \$P(T=1|S=1) = 0.64\$} \\
\texttt{    *   \$P(A=1|S=0, T=0) = 0.46\$} \\
\texttt{    *   \$P(A=1|S=0, T=1) = 0.77\$} \\
\texttt{    *   \$P(A=1|S=1, T=0) = 0.17\$} \\
\texttt{    *   \$P(A=1|S=1, T=1) = 0.43\$} \\
\texttt{    *   \$P(S=1) = 0.49\$} \\
\\ 
\texttt{**The Identification Problem:**} \\
\texttt{To calculate the causal effect \$P(A|do(S=s))\$, we typically use the adjustment formula. Since \$Tar\$ is a mediator (on the path \$S \textbackslash{}to T \textbackslash{}to A\$) and not a confounder of the \$S \textbackslash{}to A\$ relationship (it is a descendant of \$S\$), we do not adjust for it to find the *total* effect. However, we must account for the unobserved confounder \$Gender\$.} \\
\\ 
\texttt{The standard adjustment formula for an unobserved confounder \$U\$ (Gender) requires knowing \$P(S|U)\$ and \$P(A|S, U)\$. We do not have this.} \\
\texttt{Alternatively, we might try to use the observed data to estimate the effect. The observed probability of allergy for smokers is:} \\
\texttt{\$P(A=1|S=1) = P(A=1|S=1, T=0)P(T=0|S=1) +} \\
            }
        }
    }
    & & \\ \\

    \theutterance \stepcounter{utterance}  
    & & & \multicolumn{2}{p{0.3\linewidth}}{
        \cellcolor[rgb]{0.9,0.9,0.9}{
            \makecell[{{p{\linewidth}}}]{
                \texttt{\tiny{[GM$|$GM]}}
                \texttt{No} \\
            }
        }
    }
    & & \\ \\

    \theutterance \stepcounter{utterance}  
    & & & \multicolumn{2}{p{0.3\linewidth}}{
        \cellcolor[rgb]{0.9,0.9,0.9}{
            \makecell[{{p{\linewidth}}}]{
                \texttt{\tiny{[GM$|$GM]}}
                \texttt{no} \\
            }
        }
    }
    & & \\ \\

    \theutterance \stepcounter{utterance}  
    & & & \multicolumn{2}{p{0.3\linewidth}}{
        \cellcolor[rgb]{0.9,0.9,0.9}{
            \makecell[{{p{\linewidth}}}]{
                \texttt{\tiny{[GM$|$GM]}}
                \texttt{game\_result = WIN} \\
            }
        }
    }
    & & \\ \\

\end{supertabular}
}

\end{document}
